\prefacesection{Lista de Símbolos}

% \begin{center} % Remova o \centering aqui se estava usando com \begin{table}
\begin{longtable}{p{2.5cm} p{12.5cm}} % Ajuste larguras conforme necessário
    \caption[Lista de Símbolos Matemáticos e Notações]{Lista dos principais símbolos matemáticos e notações utilizados na dissertação e suas respectivas descrições. Esta tabela continua na próxima página.} \label{tab:lista_de_simbolos} \\ 

    % --- CABEÇALHO DA TABELA - Aparecerá na primeira página ---
    \toprule
    \textbf{Símbolo} & \textbf{Descrição} \\
    \midrule
    \endfirsthead

    % --- CABEÇALHO PARA PÁGINAS DE CONTINUAÇÃO ---
    \multicolumn{2}{c}%
    {{\tablename\ \thetable{} -- Continuação da página anterior}} \\
    \toprule
    \textbf{Símbolo} & \textbf{Descrição} \\
    \midrule
    \endhead

    % --- RODAPÉ PARA TODAS AS PÁGINAS, EXCETO A ÚLTIMA ---
    \midrule
    \multicolumn{2}{r}{{Continua na próxima página}} \\
    \endfoot

    % --- RODAPÉ APENAS NA ÚLTIMA PÁGINA DA TABELA ---
    \bottomrule
    \endlastfoot

    % --- CONTEÚDO DA TABELA ---
    $\alpha$  & Peso do termo de Entropia Cruzada Binária (BCE) na função de perda combinada. \\
    $\beta$   & Peso do termo de Dice Loss para a classe "câncer" na função de perda combinada. \\
    $\gamma$  & Peso do termo de Dice Loss para a classe "não câncer" na função de perda combinada. \\
    $A$       & Imagem binária de entrada para operações morfológicas. \\
    $B$       & Elemento estruturante (kernel) para operações morfológicas. \\
    $d_k$     & Dimensão dos vetores de chave (\textit{keys}) utilizada para normalização no mecanismo de atenção. \\
    $\text{Dice}(th_k)$ & Coeficiente de Dice calculado para um limiar específico $th_k$. \\
    $Dif(C_1, C_2)$ & Diferença entre dois componentes $C_1$ e $C_2$, o menor peso de aresta que os conecta. \\
    FN        & Falsos Negativos: Número de pixels incorretamente classificados como não pertencentes à classe de interesse. \\
    $\text{FN}(th_k)$ & Contagem de Falsos Negativos para um limiar específico $th_k$. \\
    FP        & Falsos Positivos: Número de pixels incorretamente classificados como pertencentes à classe de interesse. \\
    $\text{FP}(th_k)$ & Contagem de Falsos Positivos para um limiar específico $th_k$. \\
    $g_i$     & Rótulo verdadeiro (\textit{ground truth}) para o pixel $i$ (1 para classe positiva, 0 para negativa). \\
    $\hat{g}_{i}(th_k)$ & Predição binária para o pixel $i$ utilizando o limiar $th_k$. \\
    $G=(V, E)$ & Grafo que representa a imagem, com vértices $V$ (pixels) e arestas $E$. \\
    $H$       & Matriz de concentração/densidade de corantes. \\
    $I$       & Intensidade da luz transmitida (após a amostra). \\
    $I_0$     & Intensidade da luz incidente (antes da amostra). \\
    $\mathbb{I}(\cdot)$ & Função Indicadora (1 se a condição é verdadeira, 0 caso contrário). \\
    $Int(C)$ & Diferença interna de um componente $C$, o maior peso de aresta em sua árvore geradora mínima. \\
    $k$       & Parâmetro de escala que controla a preferência por componentes maiores na segmentação por grafos. \\
    $K$       & Número de limiares candidatos para otimização. \\
    $K$       & Matriz de Chaves (\textit{Keys}) no mecanismo de atenção do Transformer. \\
    $L*, a*, b*$ & Componentes do espaço de cor CIE LAB (Luminosidade, cromaticidade a, cromaticidade b). \\
    $MInt(C_1, C_2)$ & Limiar de diferença interna mínima para decidir a fusão de componentes. \\
    $N$       & Número total de pixels em uma imagem ou \textit{patch} analisado. \\
    $N_V$     & Número de \textit{patches} no conjunto de validação $V$. \\
    $OD$      & Densidade Óptica, medida logarítmica da absorbância da luz. \\
    $p_i$     & Probabilidade predita pelo modelo para o pixel $i$ pertencer à classe positiva. \\
    $P_{LAB, origem}$ & Valor de um pixel no espaço LAB da imagem original (antes da normalização). \\
    $P_{LAB, norm}$ & Valor de um pixel no espaço LAB após a normalização de cor de Reinhard. \\
    $Q$       & Matriz de Consultas (\textit{Queries}) no mecanismo de atenção do Transformer. \\
    $t$       & Limiar (threshold) calculado pelo Método de Otsu para binarização de imagem. \\
    $th_k$    & $k$-ésimo limiar candidato no processo de otimização de limiar. \\
    $th^*$    & Limiar de decisão ótimo encontrado após otimização. \\
    $Th$      & Conjunto de todos os limiares candidatos $\{th_1, ..., th_K\}$. \\
    TN        & Verdadeiros Negativos: Número de pixels corretamente classificados como não pertencentes à classe de interesse. \\
    TP        & Verdadeiros Positivos: Número de pixels corretamente classificados como pertencentes à classe de interesse. \\
    $\text{TP}(th_k)$ & Contagem de Verdadeiros Positivos para um limiar específico $th_k$. \\
    $V$       & Conjunto de validação utilizado para otimização de limiar. \\
    $V$       & Matriz de Densidade Óptica (no contexto de separação de corantes). \\
    $V$       & Matriz de Valores (\textit{Values}) no mecanismo de atenção do Transformer. \\
    $w(v_i, v_j)$ & Peso de uma aresta, medindo a dissimilaridade entre os pixels $v_i$ e $v_j$. \\
    $W$       & Matriz de vetores de cor (base de cores dos corantes). \\
    $|C|$     & Tamanho (número de pixels) de um componente $C$. \\
    $\epsilon$ & Pequeno fator de suavização (epsilon) para estabilidade numérica. \\
    $\mu_{origem}$ & Média dos valores dos pixels da imagem de origem em um canal LAB. \\
    $\sigma_{origem}$ & Desvio padrão dos valores dos pixels da imagem de origem em um canal LAB. \\
    $\mu_{ref}$ & Média de referência em um canal LAB para normalização de cor. \\
    $\sigma_{ref}$ & Desvio padrão de referência em um canal LAB para normalização de cor. \\
    $\tau(C)$ & Função de limiar adaptativa, baseada no tamanho do componente $C$. \\
    $\lambda$ & Parâmetro de regularização que controla a esparsidade na separação de corantes. \\
    $\ominus$  & Operador de Erosão morfológica. \\
    $\oplus$  & Operador de Dilatação morfológica. \\
    $\circ$   & Operador de Abertura morfológica. \\
    $\bullet$ & Operador de Fechamento morfológico. \\

\end{longtable}
